\documentclass{curriculum-vitae}
\usepackage{hyperref}
\usepackage{amsmath}
\usepackage{comment}

\author{Ian Chi}
\email{itchi@ucsc.edu}
\phone{510-304-5428}
\site{\href{https://sites.google.com/view/itchi-site/home}{sites.google.com/view/itchi-site}}




\begin{document}
\maketitle

\section{Education}
\begin{itemize}
    \item
    \itemheading{University of California, Santa Cruz (M.S.) 2025-Present}{
        Computer Science Engineering
    }
    \item
    \itemheading{University of California, Davis (B.S.) 2020-2024}{
        Physics and Mathematics (high honors), Cumulative GPA:3.859 
    }
    {}
\end{itemize}

\section{Skills Summary}
\begin{itemize}
    \item
    % \begin{itemshortlist}
        \textbf{Programming}:

        \begin{itemize}
            \item Experienced: Python, Mathematica 
    
            \item Other: C, C\#, Java, Carnap, Rocq
        \end{itemize}
    % \end{itemshortlist}

    \item
    \begin{itemshortlist}{Languages}
        English, Chinese (Mandarin)
    \end{itemshortlist}

    \item
    \begin{itemshortlist}{Others}
        \LaTeX, LT Spice, Lean
    \end{itemshortlist}
\end{itemize}

\section{Publications \& Research}
\begin{itemize}
    \item\textbf{Publications \& Preprints}:
    \begin{enumerate}
        \item\label{pub: pub1} I. Chi, T. Tan, M. Fraas. \textbf{``Poincar\'e disk as a model of squeezed states of a harmonic oscillator''}. \\
        %Squeezed states under specific Quadratic Hamiltonians have nice geometric properties when mapped to Poincare Disk. This paper goes over properties and applications of said properties in quantum control theory. \\
        Paper Link: \hyperlink{https://pubs.aip.org/aip/jmp/article-abstract/66/6/062101/3349274}{https://pubs.aip.org/aip/jmp/article-abstract/66/6/062101/3349274}

        Honors Thesis Link: \hyperlink{https://tinyurl.com/2pj46txh}{www.math.ucdavis.edu/\~webfiles/undergrad\_thesis/202403\_Ian\_Chi\_Fraas.pdf}

       

    \end{enumerate}
    \item \textbf{Research}
    \begin{enumerate}
        \item \textbf{Simplifying and Generalizing Polylogarithmic-Time Longest Increasing Subsequence (LIS) Algorithms}
        
        Ongoing work re-examining the Saks \& Seshadhri (2010) foundational sublinear-time algorithm. Deconstructing its proofs to improve conceptual clarity, generalize the framework to other DP problems, and optimize constant factors for sublinear LIS estimation.
        \item \label{pub: pub2} I. Chi, R. M. Morin, R. D. Mecholsky, N. A. Mecholsky. \textbf{``Derivation of Free Flight Distance Distribution with Application to Drift Velocity Calculations''}.
         
        Derived probability distributions for free-flight distances and applied statistical models to electron drift velocity calculations within the Drude framework.
        (Funded by NSF Grant PHY-1852006; presented poster available \href{https://drive.google.com/file/d/1rUFjoE8ORksClzMPz6evIloqEMJk0EOd/view}{\emph{here}}).
    \item \textbf{Laboratory Researcher, DUNE \& ANNIE Neutrino Experiments} (UC Davis / Lab) \\
    Worked with Dr. Robert Svoboda on developing and testing Water Based Slow Scintillators (WBSS) to isolate Cherenkov radiation from scintillation signals. Designed and constructed a double coincidence light yield test bench and performed data analysis on scintillation decay times.
    \end{enumerate}
    
\end{itemize}


\section{Selected Coursework (UC Santa Cruz, M.S.)}
\begin{itemize}
    \item\textbf{Computer Science Engineering}
    \begin{itemize}
        \item Approximation Algorithms, Complexity Theory, Analysis of Algorithms, Programming Languages, Machine Learning Theory, Modern Algorithms
    \end{itemize}


\end{itemize}

\section{Selected Coursework (UC Davis, B.S.)}
\begin{itemize}
    \item\textbf{Math courses}
    \begin{itemize}
        \item Real \& Complex Analysis, Modern Algebra, Topology, Functional Analysis (Grad), Algebraic Topology (Grad), Probability.
    \end{itemize}


    \item\textbf{Physics courses}:
    \begin{itemize}
        \item Classical Mechanics,
        Electricity and Magnetism,
        Quantum Mechanics,
        Electronics Lab,
        Particle Physics,
        General Relativity,
        Thermal and Statistical Mechanics,
        Computational and Mathematical,

    \end{itemize}

    \item\textbf{Philosophy courses}:
    \begin{itemize}
        \item Symbolic Logic,
        Philosophy of Mathematics, 
        Philosophy of Science,
        Philosophy of Sexuality
    \end{itemize}
\end{itemize}




\section*{Awards}
\begin{itemize}
    \item \textbf{Citation for Outstanding Performance, UC Davis Mathematics Department} 
    \item \textbf{Citation for Outstanding Performance, UC Davis Physics Department} 
    \item \textbf{GMTK Game Jam 2023 -- Honorable Mention (Top Tier among 6,000+ Entrants)}\\
    Served as Lead Designer for an independent team developing \emph{SPARED!}, a rapid-prototype puzzle game built under a strict 48-hour deadline. Designed core game mechanics, balancing systems, and player onboarding loops, earning an Honorable Mention out of about 7,000 global entries. 
    \newline Portfolio \& Playable Builds: \url{https://splake975.itch.io/}
    % \item \textbf{Dean's List, School of Letters and Sciences} -- Fall 2020 \& Spring 2021
\end{itemize}

\section{Extracurricular Activities \& Experiences}
\begin{itemize}
    \item \textbf{Host of the UCSC TCS Seminar} \emph{\href{https://cstheory.ucsc.edu/theory-seminar/}{https://cstheory.ucsc.edu/theory-seminar/}} (Fall 2026) \\
    Organizing and hosting weekly theory seminars, facilitating discussions on cutting-edge research in theoretical computer science.

    \item \textbf{Applied Probability \& Software Project: Advantage Play Modeling} -- (Feb 2024 -- Dec 2025)\\
    Designed and implemented software to analyze statistical edges and manage game states for advantage play. Generated \$600k+ in revenue with a 7\%+ return over the period. Project concluded following regulatory shifts in California.

\item \textbf{Nuclear Science and Security Consortium (NSSC) Fellow} (Spring 2023 -- Summer 2024) \\
    Selected for a competitive multidisciplinary fellowship focused on advanced nuclear security and instrumentation.

\item \textbf{William Lowell Putnam Mathematical Competition} -- 14 points (December 2023) \\
    Achieved a competitive score in North America's premier undergraduate mathematics competition, demonstrating advanced problem-solving skills under tight time constraints.


\item \textbf{Directed Reading Program -- Lie Groups and Lie Algebras} (Spring 2023) \\
    Departmental reading program under Michael Ragone (PhD Candidate). Studied differential geometry and representation theory via Tu's \emph{Introduction to Manifolds} and Hall's \emph{Lie Groups, Lie Algebras and Representations}.

\item \textbf{Directed Reading Program -- Differential Topology and Geometry} (Fall 2022 -- Winter 2023) \\
    Studied differential forms, co/tangent spaces, and manifolds. Concluded with a departmental presentation under faculty mentorship.

\item \textbf{Research on Rails Research Trip (Topological Quantum Field Theory)} \\
    Led by Dr. Andrew Waldron; worked with PhD candidate Petr Vlachopulos studying Michael Atiyah's foundational paper on TQFTs, focusing on the categorical structures of cobordisms and linear maps.

% \item \textbf{2023 Physics Research Experience for Undergraduates (REU)} \\
%     Conducted research at Catholic University of America under Dr. Nicholas A. Mecholsky and Rachel M. Morin. (See poster reference \ref{pub: pub2}).
    
\end{itemize}









\end{document}